\section{A polynomial error unless a short cycle creates a long detour}
\label{sec:second-order-detours}

The exponential cut estimate also sharpens a structural alternative in the unrestricted proof.
A polynomial excess above one half forces a short cycle whose deletion greatly separates two of its neighbors.
This gives a polynomial cycle-count error for fixed cubic vertex-expanders.
It does not give that error for arbitrary graphs.

All logarithms below have base two unless written as $\ln$.
Fix the absolute constants
\[
 A=6e^{16},\qquad t=\lceil1000A\rceil,
 \qquad \alpha=\frac1{10^6t2^t}>0.
\]
Their sizes are not optimized.

\begin{lemma}[A terminal forest meets every multiple-link partner]
\label{lem:multiple-link-terminal-forest}
Let $C$ be a cycle of a connected subcubic graph $\Gamma$.
For each component $H$ of $\Gamma-V(C)$, choose a tree $T_H\subseteq H$ connecting all its neighbors of $C$.
Include a single vertex when there is only one such neighbor.
Let $T$ be the union of these trees.
Every cycle $C'$ disjoint from $C$ with at least two links to $C$ meets $T$.

If every pair of neighbors of $C$ in the same component has distance at most $\Delta$ there, one may choose
\[
 |V(T)|\le |C|(\Delta+1).
\]
\end{lemma}

\begin{proof}
Suppose $C'$ avoids $T$.
It lies in one component $H$ of $\Gamma-V(C)$.
It contains no neighbor of $C$, so no edge directly joins the two cycles.
After deleting $C'$, the tree $T_H$ still connects all neighbors of $C$ in $H$.
Thus at most one component of $\Gamma-V(C)-V(C')$ can be joined to both cycles.
This gives at most one link.

For the size bound, fix one neighbor in each component.
Join every other neighbor to it by a path of length at most $\Delta$.
Take a spanning tree of the resulting connected union.
The total number of terminal incidences is at most $|\delta V(C)|\le|C|$.
Summing the path-length bounds proves the assertion.
\end{proof}

\begin{theorem}[A polynomial-gap detour obstruction]
\label{thm:polynomial-gap-detour}
Let $\Gamma$ be a connected multigraph of rank $r$, with all degrees two or three.
Let $P\subseteq V(\Gamma)$ contain every degree-two vertex, and suppose
\[
 |P|\le\frac{2r}{\log^2r},\qquad
 \operatorname{comp}(\Gamma[P])\le\frac{\log r}{100}.
\]
Let $\mathcal F$ mean that a uniform even subgraph is a forest on $\Gamma-P$.
For all sufficiently large $r$, the inequality
\[
 \Pr(\mathcal F)>\frac12+r^{-\alpha}
\]
implies the following conclusion.
There is a cycle $C\subseteq\Gamma-P$ of length at most $A\log r$ and a component $H$ of $\Gamma-V(C)$ containing two neighbors $u,v$ of $C$ with
\[
 \operatorname{dist}_H(u,v)>r^{1/25}.
\]
The threshold is absolute.
\end{theorem}

\begin{proof}
Put $x=\log r$, $y=\log x$, $D=\lfloor x/10\rfloor$, $s=\lfloor(D-1)/2\rfloor$, and $L=\lfloor Ax\rfloor$.
Assume the displayed probability bound and suppose the detour conclusion fails.

Choose a maximal disjoint family $\mathcal Q$ of cycles of length at most $D$ in $\Gamma-P$.
Its cycle regions have cuts of size at most $D$.
Theorem~\ref{thm:sharper-exponential-cuts} gives
\[
 |\mathcal Q|<8D2^D\ln(2r^\alpha).
\]
Let $S$ be their vertex union.
Then
\[
 |S|=O(x^3r^{1/10})=o(r/x).
\]
The graph $J=\Gamma-(P\cup S)$ is simple and has girth greater than $D$.
Its vertex count $v$ and degree deficit $d_J$ satisfy
\[
 v=2r+O(r/x^2+|S|),\qquad
 d_J\le3(|P|+|S|),\qquad Ld_J=o(v).
\]

Let $\mathcal C$ contain its cycles of length at most $L$, and write
\[
 a_C=2^{-|C|},\qquad \mu=\sum_{C\in\mathcal C}a_C.
\]
The paired-power spectral lemma and the repeated-walk estimate give
\[
 \mu\ge\frac12\ln(A/6)-o(1)=8-o(1)\ge7.
\]
For every vertex $u$ of $J$, the manuscript's walk bound gives
\begin{equation}\label{eq:detour-local-cycle-mass}
 \sum_{C\in\mathcal C:u\in V(C)}a_C\le L2^{-s}.
\end{equation}

We next check the exceptional-pair estimate with the smaller error margin.
Use the fixed integer $t$ in the link decomposition.
For two disjoint cycles $C,C'\in\mathcal C$, the cyclic complementary regions with cuts at most $t$ number fewer than
\[
 8t2^t\ln(2r^\alpha)
 =\frac{8\ln2}{10^6}x+O_t(1).
\]
The other terms in the link decomposition are unchanged.
For sufficiently large $r$, one has $2t+D/2\le D$, and hence
\[
 \ell(C,C')\le \operatorname{comp}(\Gamma[P])
 +\frac{2L}{t}+8t2^t\ln(2r^\alpha)+(8L/D)^2
 \le x/40.
\]
Indeed, the four terms are at most $x/100,x/500,x/1000,x/1000$ eventually.
The exceptional-partner lemma therefore bounds the normalized exceptional contribution for each fixed $C$ by
\[
 2^{x/40}L^3 2^{-s}=o(1).
\]

For $C\in\mathcal C$, let $w_2(C)$ be the total weight of its disjoint partners with at least two links.
The assumed absence of the detour conclusion and Lemma~\ref{lem:multiple-link-terminal-forest} give a hitting forest with at most $L(r^{1/25}+1)$ vertices.
Using \eqref{eq:detour-local-cycle-mass}, including only its vertices in $J$, gives
\[
 w_2(C)\le L^2(r^{1/25}+1)2^{-s}=o(1),
\]
uniformly over $C$.

Use the manuscript's normalized component indicators $X_C$ and let $Z=\sum_CX_C$.
Then $\mathbb EZ=\mu$ and $\mathbb EX_C^2\le a_C$.
Intersecting cycles contribute zero.
One-link pairs contribute exactly $a_Ca_{C'}$.
Two-link pairs contribute twice that amount, and exceptional pairs contribute $o(\mu)$ in total.
Consequently,
\[
 \mathbb EZ^2
 \le\mu^2+\mu+\sum_Ca_Cw_2(C)+o(\mu)
 =\mu^2+(1+o(1))\mu.
\]
Since $\mathcal F$ implies $Z=0$, Cauchy--Schwarz gives
\[
 \Pr(\mathcal F)\le\Pr(Z=0)
 \le\frac{1+o(1)}{\mu+1+o(1)}\le\frac18+o(1).
\]
This contradicts the assumed probability above one half.
\end{proof}

\begin{lemma}[Diameter after deleting vertices from an expander]
\label{lem:expander-deletion-diameter}
Suppose an $n$-vertex graph $Q$ satisfies
\[
 |N_Q(U)\setminus U|\ge h|U|
 \qquad(0<|U|\le n/2)
\]
for some $0<h\le1$.
For any set $Z$ of $z$ vertices, every component of $Q-Z$ has diameter at most
\[
 \frac{4z+6\ln n+4}{h}.
\]
\end{lemma}

\begin{proof}
A component $K$ of size at most $n/2$ has $h|K|\le z$, so its diameter is at most $z/h$.
There is at most one larger component.
If $z/h\ge n/4$, the claimed bound follows from the trivial diameter bound $n-1$.
Assume $z/h<n/4$.
In that component, a ball $B$ of size at most $n/2$ has at least $h|B|-z$ external neighbors within the component.
Connectivity grows a ball to size at least $\max\{1,2z/h\}$ in at most $2z/h+1$ steps.
After that point, every step multiplies its size by at least $1+h/2$, until it exceeds $n/2$.
Since $\ln(1+h/2)\ge h/3$, this takes at most $3\ln n/h+1$ further steps.
Two such balls intersect.
Their combined radius is at most $(4z+6\ln n+4)/h$, as required.
\end{proof}

\begin{corollary}[The forest bound for fixed expanders]
\label{cor:expander-polynomial-forest}
Under the hypotheses of Theorem~\ref{thm:polynomial-gap-detour}, suppose additionally that $\Gamma$ has vertex expansion at least a fixed $h>0$.
Then, for sufficiently large $r$ depending on $h$,
\[
 \Pr(\mathcal F)\le\frac12+r^{-\alpha}.
\]
\end{corollary}

\begin{proof}
The vertex count is $2r+o(r)$.
Deleting any cycle of length at most $A\log r$ leaves components of diameter $O_h(\log r)$.
For large $r$, this is smaller than $r^{1/25}$.
Apply Theorem~\ref{thm:polynomial-gap-detour}.
\end{proof}

\begin{theorem}[A polynomial cycle error for cubic expanders]
\label{thm:expander-polynomial-cycles}
Fix $h>0$.
For every sufficiently large connected simple cubic graph $G$ with vertex expansion at least $h$ and rank $r$,
\[
 c(G)\le2^{r-1}+2\frac{2^r}{r^\alpha}.
\]
\end{theorem}

\begin{proof}
Choose a spanning tree of $G$.
Its $r$ cotree edges contain a matching of size at least $r/5$, by greedy selection in a graph of maximum degree three.
Choose a matching $D$ of $k=\lceil\alpha\log_2r\rceil$ cotree edges.

For each $uv\in D$, replace that edge by
\[
 u\!-!u'\!-!v'\!-!v,
 \qquad u'\!-!w\!-!v',
\]
using three new vertices.
The middle edge $u'v'$ and the two edges through $w$ form a triangle.
Call the new graph $\Gamma$ and let $P$ be the union of these triangle vertices.
Then $\beta(\Gamma)=R=r+k$.
Every unmarked vertex has degree three, every marked vertex has degree two or three, and
\[
 |P|=3k,\qquad \operatorname{comp}(\Gamma[P])=k.
\]
The matching condition ensures that the gadgets have distinct endpoints.

We check expansion explicitly.
Partition $V(\Gamma)$ into connected branch sets indexed by $V(G)$, each of size at most three.
Assign $u',w$ to $u$ and $v'$ to $v$ for each gadget.
Every original edge has a representative between the corresponding branch sets.

For $U\subseteq V(\Gamma)$ with $|U|\le|V(\Gamma)|/2$, let $B$ count partially occupied branch sets.
Each such branch set supplies a distinct external neighbor of $U$, so $|N_\Gamma(U)\setminus U|\ge B$.
The numbers of fully occupied and fully unoccupied branch sets are each at least $|U|/3-B$.
The original expansion and maximum degree three therefore give
\[
 |\delta_\Gamma U|\ge h|U|/3-(h+3)B.
\]
If $B\ge h|U|/[6(h+3)]$, the first bound suffices.
Otherwise the edge boundary has size at least $h|U|/6$.
Since $\Gamma$ has maximum degree three, its external vertex boundary then has size at least $h|U|/18$.
Thus $\Gamma$ has vertex expansion at least $h/[6(h+3)]$.

The size and component conditions of Corollary~\ref{cor:expander-polynomial-forest} hold for large $r$.
It gives $\Pr(\mathcal F)\le1/2+R^{-\alpha}$.
Lift every original cycle meeting $D$ through the three-edge gadget paths.
Toggle any subset of the $k$ triangles.
The resulting $2^k$ even subgraphs are distinct, and their restrictions to $\Gamma-P$ are forests.
The new edges through each $w$ determine the toggled subset, so the full map is injective.
Hence the cycles meeting $D$ number at most $\Pr(\mathcal F)2^r$.

At most $2^{r-k}$ cycles avoid $D$, since the chosen cotree coordinates are independent.
Finally, $R^{-\alpha}\le r^{-\alpha}$ and $2^{-k}\le r^{-\alpha}$.
Combining these estimates proves the theorem.
\end{proof}

\begin{corollary}[Expansion may decay polynomially]
\label{cor:weak-expander-polynomial-cycles}
For all sufficiently large connected simple cubic graphs $G$ of rank $r$ with vertex expansion at least $r^{-1/50}$,
\[
 c(G)\le2^{r-1}+2\frac{2^r}{r^\alpha}.
\]
The threshold is absolute.
\end{corollary}

\begin{proof}
Apply the same triangle-gadget construction.
Its rank is $R=r+O(\log r)$, and its vertex expansion is at least $r^{-1/50}/24$.
Indeed, the previous comparison gives $h/[6(h+3)]\ge h/24$ for $0<h\le1$.
Deleting a cycle of length at most $A\log R$ therefore leaves components of diameter
\[
 O(r^{1/50}\log r)<R^{1/25}
\]
for all sufficiently large $r$.
The constant in this estimate is absolute, by Lemma~\ref{lem:expander-deletion-diameter}.
The detour theorem now gives the same forest estimate.
The triangle-flip injection and the witness-avoidance bound complete the proof as before.
\end{proof}

The obstruction theorem does not yet control the case of long terminal detours.
Such detours suggest distance-layer decompositions.
A proof must still combine the resulting localized restrictions without losing the density gap.
